\section{Results on \mgp{}}
\label{sec:results}

All results below are recomputed from released per-item predictions. We lead
with validity checks because they determine how the policy scores may be
interpreted.

\subsection{MGP-6 v4.3 final hidden results: pending}

The v4.3 release is integrated with public train/development data and a private
hidden scoring boundary, but final hidden predictions have not yet been
generated under a frozen configuration. Accordingly, this manuscript does not
report v4.3 accuracy, macro-F1, confidence intervals, confusion matrices,
ablations, cost, latency, or paired comparisons. The final-results block will
be inserted only after predictions are finalized and hashed before private
scoring.

\begin{center}
\fbox{\parbox{0.88\linewidth}{\centering
\textbf{Placeholder:} MGP-6 v4.3 final hidden system comparison,\
component-cluster confidence intervals, confusion matrix, ablations, and
cost/latency table.}}
\end{center}

\subsection{The shortcut check fails}

The conflict-type-only lookup fitted on train+dev reaches $1.0000$ accuracy and
macro-F1 on test; the coarser six-way type lookup reaches $0.9667$ accuracy and
$0.9471$ macro-F1. Neither sees the dialogue or memory state. Thus the held-out
item split does not test whether a system reasons about evidence quality,
provenance, or temporal scope: scenario recognition is sufficient. This failed
check dominates every policy comparison below.

\begin{table}[t]
\centering
\small
\caption{Shortcut probes fitted on train+dev and evaluated on the held-out test
split. These probes never read the dialogue or memory state.}
\label{tab:shortcut}
\begin{tabular}{@{}lrrr@{}}
\toprule
Probe & Test accuracy & Macro-F1 & Unseen values \\
\midrule
Six-way type only & 0.9667 & 0.9471 & 0 \\
Conflict type only & 1.0000 & 1.0000 & 0 \\
\bottomrule
\end{tabular}
\end{table}


\subsection{Held-out policy comparison}

% Source: per-item predictions under artifacts/mgp6_v2; test split only.
\begin{table}[t]
\centering
\small
\caption{Held-out test results ($n{=}90$). Family CIs resample the 48
template families and report family-macro accuracy. Action false-commit is over 33
gold-\abstain{} items; missed-update is over 45 gold-\supersede{} items. The
shortcut result in Table~\ref{tab:shortcut} limits interpretation of all rows.}
\label{tab:main}
\begin{tabular}{@{}lrrrrr@{}}
\toprule
Policy & Accuracy & Family 95\% CI & Macro-F1 & Action false-commit & Missed-update \\
\midrule
Reference policy & 0.8889 & [0.8125, 0.9688] & 0.9145 & 0.0000 & 0.2222 \\
Structured LLM & 0.9000 & [0.8333, 0.9688] & 0.9040 & 0.0606 & 0.1556 \\
Unstructured zero-shot LLM & 0.6778 & [0.5625, 0.7917] & 0.6542 & 0.5758 & 0.2222 \\
\bottomrule
\end{tabular}
\end{table}


On test, the reference policy reaches $0.8889$ accuracy/$0.9145$ macro-F1 and
the structured LLM reaches $0.9000$/$0.9040$. The reference-minus-LLM risk
difference is $-0.0111$; the paired template-family cluster interval
$[-0.094,0.063]$ includes zero (exact item-level McNemar:
$b=4,c=5,p=1.0$). These data support neither superiority nor non-inferiority.
The unstructured zero-shot policy reaches $0.6778$/$0.6542$, but that gap is
itself evaluated on the shortcut-leaking pilot and is not a general
scaffold-effect estimate.

\subsection{Safety--responsiveness trade-off}

The reference policy makes no action false commits on the 33 test abstention items, but
the exact one-sided 95\% upper bound is $0.0868$, not zero risk. It misses
$10/45$ required updates ($0.2222$) and unnecessarily abstains on $10/57$
non-abstention items ($0.1754$). The structured LLM action false-commit rate is
$2/33=0.0606$, while its missed-update rate is $7/45=0.1556$. Consequently the
reference policy has lower expected loss under the commit-averse matrix
($0.111$ vs.\ $0.189$), but higher loss under the update-responsive matrix
($0.556$ vs.\ $0.433$). Safety is therefore cost-conditional, not an absolute
ordering. Top-label ECE is $0.156$ for the reference confidence and $0.302$ for
the LLM confidence; both are descriptive because neither confidence was
calibrated on an independent set.

\subsection{State and component diagnostics}

Full-corpus ablations are retained only as within-generator mechanism checks.
Removing abstention raises invalid-commit from $0.0$ to $1.0$; removing
provenance yields $0.606$; removing typing drives \keep{} F1 to $0.0$. Removing
lineage leaves action accuracy unchanged at $0.9222$ but drops lineage integrity
$0.8556\to0.3689$. The last result demonstrates the central measurement point:
action accuracy alone cannot reveal a corrupted transition history.

% tables/ablation.tex
% All values from artifacts/results_ledger.json v1.1 -> mgp6_v2.results[*].
% The later shortcut audit makes these within-generator diagnostics only.
%   ablation/full         : action_accuracy 0.9222, macro_f1 0.9399   (= rule controller)
%   ablation/no_abstain   : 0.5556 / 0.4278 / invalid_commit 1.0
%   ablation/no_provenance: 0.7000 / 0.7428 / invalid_commit 0.606
%   ablation/no_typing    : 0.7889 / 0.5680 / keep_f1 0.0
%   ablation/no_lineage   : 0.9222 / 0.9399 / lineage 0.3689
% Full-system lineage integrity 0.8556 from ref_5systems/controller_rule.lineage.
\begin{table}[t]
\centering
\small
\caption{Within-generator component ablations of the reference policy on \mgp{}
($n{=}180$). Each removal produces a distinct diagnostic signature.
Invalid-commit is over gold-\abstain{} items; lineage integrity is over
gold-\supersede{} items. Values transcribed from
results-ledger version 1.1; these are not independent effect estimates.}
\label{tab:ablation}
\begin{tabular}{@{}l r r l@{}}
\toprule
Configuration & Action acc. & Macro-F1 & Induced failure signature \\
\midrule
Full system        & 0.9222 & 0.9399 & invalid-commit 0.0; lineage integrity 0.8556 \\
$-$ abstention      & 0.5556 & 0.4278 & invalid-commit $0.0 \to 1.0$ \\
$-$ provenance      & 0.7000 & 0.7428 & invalid-commit $0.0 \to 0.606$ \\
$-$ typing          & 0.7889 & 0.5680 & \keep{} F1 $\to 0.0$ \\
$-$ lineage         & 0.9222 & 0.9399 & action unchanged; lineage integrity $0.8556 \to 0.3689$ \\
\bottomrule
\end{tabular}
\end{table}


The $174/180$ six-shot run remains incomplete and is excluded from all
comparisons. Mapped external-system action scores and non-homogeneous downstream
runs remain in the artifact audit only; neither supports a main-paper capability
claim.
